\subsection{The uniform thickening argument and a precise corrected theorem}

The next result supplies the conditional argument previously recorded only
in our Artin--Rees research note. We give the module statement first,
since its uniformity is the reason that condition 8 can be used on one
finite thickening, rather than on infinitely many different thickenings.

\begin{lemma}[A uniform finite-module test]\label{A:uniform-AR}
\status{Proved}
Let $R$ be noetherian, let $I\subset R$ be an ideal, and let $C$ be a
finite $R$-module. There is $N\geq1$ with the following property. For
any flat noetherian local $R$-algebra $R'$ with $IR'$ in its maximal
ideal, any finite $R'$-module $D$, and any homomorphism
$\phi:D\to C\otimes_RR'$, if $\phi$ modulo $I^N$ is an isomorphism,
then $\phi$ is an isomorphism. The same $N$ works with any direct
summand of $C\otimes_RR'$ as target.
\end{lemma}
\begin{proof}
Choose a presentation $0\to K\to R^d\to C\to0$.
Artin--Rees, in the form of Stacks, Tag 00IN
\cite{Stacks00IN}, gives $N$ such that
$K\cap I^NR^d\subset IK$. Flat base change preserves the kernel
and this intersection. We may therefore work over $R'$ and omit
primes in the notation. Nakayama makes $\phi$ surjective. Lift the
chosen $d$ generators of $C$ to $D$. They generate $D$ modulo
$I^N$, hence generate $D$. If $K'$ is the kernel of the resulting
$R^d\to D$, then $K'\subset K$. The isomorphism modulo $I^N$ says
$K\subset K'+I^NR^d$. Subtracting an element of $K'$ shows
\[
 K\subset K'+(K\cap I^NR^d)\subset K'+IK.
\]
Nakayama applied to the finite module $K/K'$ gives $K=K'$.
For a direct summand $C_1$, add the complementary summand $C_2$
to $D$ and apply the result to $D\oplus C_2\to C_1\oplus C_2$.
\end{proof}

\begin{lemma}[Finite-fibre etale morphisms over a noetherian base]
\label{A:etale-finitefibres-qc}
\status{Known facts, with the needed deduction proved}
If $P\to T$ is separated and etale, $T$ is noetherian, and all its
fibres are finite, then $P\to T$ is quasi-compact.
\end{lemma}
\begin{proof}
The degree of an etale fibre is the number of its geometric points.
By Stacks, Tag 086R \cite{Stacks086R}, the locus where this degree
is at least $d$ is open. Every point lies on an irreducible component
of $T$, and every open containing that point contains the generic
point of this component. The finite degrees at the finitely many
generic points therefore give a uniform bound on all fibre degrees.
The bound is unchanged by field extension. Stacks, Tag 07RY
\cite{Stacks07RY}, now gives a finite locally closed stratification
of $T$ over which $P$ is finite locally free. Each inverse image
stratum is quasi-compact. Their finite union is $P$, which is
quasi-compact. Apply the same argument over every affine open of
$T$ to obtain the assertion about the morphism.
\end{proof}

\begin{theorem}[Condition 8 after actual local representability]
\label{A:local-models-suffice}
\status{Proved under the actual local-representability hypothesis}
Let $S$ be locally noetherian. Suppose $F$ satisfies conditions
1--4, 6 and 8, with finite-type tests in 8. Suppose also that, for
every $s\in S$, $F_{S_s}$ is represented by a separated locally
quasi-finite scheme, where $S_s=\Spec\OO_{S,s}$.
Then $F$ is represented by a separated locally quasi-finite
$S$-scheme. In particular every local-modularity locus $M(T,\eta)$
on a locally finite-type $S$-scheme is open, and hence locally
constructible. The quasi-modular charts and their gluing in
Theorem~\ref{A:local-to-global} consequently apply.
\end{theorem}
\begin{proof}
Define $E(T,\eta)$ to be the locus where the morphism from $T$
to the relevant actual local representing scheme is etale.
This is independent of the chosen local representative by Yoneda,
and it is stable under generization. The same local representations
exist after a noetherian base change: at a point of the new base,
its local scheme maps to $S_s$, and one takes the base change of
the representing scheme there.

\emph{Constructibility of $E$.} Work on affine finite-type pieces.
Let $C\subset T$ be irreducible and closed, with generic point $c$.
If $c\notin E$, generization stability gives $C\cap E=\varnothing$.
Suppose $c\in E$, and write $s$ for its image. Then $T\to S$ is
quasi-finite at $c$, since its map to the locally quasi-finite local
model is etale there. After shrinking $T$ around $c$ and making an
etale neighbourhood of $s$ in $S$, we can isolate a finite branch
$T\to\Spec R$ through a point above $c$. We can moreover arrange
that the map of local models at that point is a local isomorphism.
Here are the precise facts behind these two replacements. Over the
strict henselization of $\OO_{S,s}$, the branches through the closed
fibre are finite local schemes, by Stacks, Tags 04GH and 06DD
\cite{StacksHenselianBranches}. Their residue fields are purely
inseparable over a separably closed field. An etale map between
two such local branches has equal residue fields, hence is an
isomorphism: it is finite etale between henselian local rings,
and its residue degree is one. The finite branch, its decomposition,
and this finite-presentation isomorphism descend to a finite stage
of the strict-henselization system, and then to an etale neighbourhood
of $s$. This replacement is etale on $T$, so relative opens in the
pulled-back stratum have open images in $C$.

Let $I$ be the prime ideal defining the closure of the image of
$c$ in this new affine base. Write $T=\Spec C_R$, with $C_R$ finite
over $R$, and choose $N$ from Lemma~\ref{A:uniform-AR} for the
underlying $R$-module of $C_R$. Inside
$T\times_R\Spec(R/I^N)$ remove all irreducible components other
than the one whose generic point is the chosen $c$, by deleting
their union. The resulting open $Z$ is irreducible and contains
$c$. It is a finite-type scheme over the irreducible noetherian
scheme $\Spec(R/I^N)$. Its generic local map to $F_{R/I^N}$ is an
isomorphism by the preceding branch construction, so it is actually
modular. Condition 8 gives a nonempty open $W\subset Z$ on which
that map is modular (the weaker etale alternative is enough).

We verify that every point $w\in W$ belongs to the original $E$.
Let $q$ be its base image and pass to the strict henselization
$R'$ of $R_q$. The finite algebra $C_R\otimes_RR'$ decomposes as
a product of finite local algebras. The branch $C_w$ through a
chosen geometric point above $w$ is one factor. A finite local
branch $D$ of the actual representing scheme exists by Tag 04GH;
the morphism gives $D\to C_w$. Modulo $I^N$ this is an etale map
of finite local henselian algebras, with equal residue fields by
Tag 06DD, and so is an isomorphism. The branch lies entirely in
$W$ after reduction: an open containing the closed point of a local
scheme is the whole scheme. Lemma~\ref{A:uniform-AR}, including
its direct-summand assertion, therefore makes $D\to C_w$ itself
an isomorphism. Faithfully flat etale descent of etaleness proves
$w\in E$ before strict henselization.

The component through $c$ is, near its generic point, the pulled-back
stratum $C$: for a finite morphism, primes above $I$ are incomparable
and are the generic points of the inverse image of $V(I)$. Removing
the other components makes this identification on an open of that
stratum. Thus the etale image of $W$ supplies a nonempty relative
open of the original $C$ contained in $E$. Noetherian induction
now proves constructibility of $E$. A constructible subset stable
under generization in a noetherian space is open. This proves
openness of $E$ on every finite-type test.

\emph{Etale charts.} Fix an affine finite-type open $V_s$ of a
local representing scheme. Finite presentation spreads it to an
affine finite-type $V$ over an affine neighbourhood of $s$ in $S$,
and condition 3 spreads its universal element to $\xi\in F(V)$.
The open locus $E(V,\xi)$ contains $V_{S_s}$. Its closed complement
is finite type and has empty localization at $s$, so finitely many
denominators outside $s$ kill it. Shrink the base accordingly.
We obtain a chart everywhere etale in the local-model sense.

For a noetherian $S$-scheme $T$ and $\eta\in F(T)$, put
$P=h_V\times_Fh_T$. It satisfies conditions 1--4 relative to $T$,
since those conditions are preserved by fibre products. Its
Artinian nilpotent lifting maps are bijective, as can be checked
in the actual local representing schemes, where they are the
lifting maps of the etale chart. Valuative uniqueness holds because
$V\to S$ is separated. The proved etale corollary,
Proposition~\ref{A:etale-corollary}, represents $P$ by a separated
etale $T$-scheme. For an arbitrary affine $T$, condition 3 descends
$\eta$ to a finite-type $T_i$ over an affine piece of $S$; the
fibre product over $T$ is the base change of the one over $T_i$.
Zariski gluing handles arbitrary $T$. Hence $h_V\to F$ is a
representable etale morphism on all test schemes.

These charts cover all field-valued points, by their construction
from all the local representing schemes. Their represented etale
base changes therefore form surjective etale covers. The scheme
$U=\coprod V$ and its scheme relation $R=U\times_FU$ give an
etale presentation. Stacks, Theorem 65.10.5, Tag 02WW
\cite{Stacks02WW}, represents $F$ by an algebraic space $X$.
Local quasi-finiteness follows on the charts.

\emph{Separation and schemeness.} For two affine charts $V_i,V_j$,
$R_{ij}=V_i\times_XV_j\to V_i$ is separated and etale: its map
into $V_i\times_SV_j$ is a monomorphism, and the second chart is
affine over the chosen affine base. Its fibres are finite. Indeed,
over the local representing scheme at a point of $S$ they are
fibres of an etale map with noetherian source $V_j$, and hence
quasi-compact zero-dimensional fibres. Lemma
\ref{A:etale-finitefibres-qc} makes $R_{ij}$ quasi-compact.
Thus $X\to S$ is quasi-separated. Valuative uniqueness and
Stacks, Tag 03KV \cite{Stacks03KV}, make it separated. Equivalently,
if condition 6 is tested only on discrete valuation rings, apply
the noetherian finite-type valuative criterion to the diagonal
on the affine chart pairs. Stacks, Proposition 67.50.2, in Tag 0417
\cite{Stacks0417}, now shows that the separated locally quasi-finite
space $X$ is a scheme.

Finally, for a locally finite-type morphism of locally noetherian
schemes, an isomorphism on the two local rings at a point spreads
to an isomorphism on open neighbourhoods. One first restricts to
the etale locus and then spreads the inverse local-ring map by
finite presentation. Consequently the local-isomorphism locus
$M(T,\eta)$ is open. This proves the asserted L1 conclusion under
exactly the local-representability premise stated above.
\end{proof}

\begin{theorem}[The pro-modular correction of condition 7]
\label{A:pro7-representability}
\status{Proved under the explicitly strengthened condition $7_{\mathrm{pro}}$}
Replace the actual-modularity premise of condition 7 by
pro-modularity at the closed point: for every complete noetherian
local base $A$, every finite local $A$-algebra $B$, and every
$\xi\in F(B)$ which represents the corresponding local Artinian
deformation functor, require that $\Spec B\to F_A$ be locally
modular, or merely etale, at every point of $\Spec B$.
Call this condition $7_{\mathrm{pro}}$.
Then conditions 1--6, $7_{\mathrm{pro}}$ and 8, with finite-type
tests in 8, are equivalent to representability by a separated
locally quasi-finite scheme over the locally noetherian $S$.
\end{theorem}
\begin{proof}
We prove sufficiency by induction on the dimension of noetherian
bases, simultaneously for all their restrictions of $F$.
The dimension-zero case is Lemma~\ref{A:artinian-base}, with
Zariski gluing of the Artinian components.

Let $(A,\mathfrak m)$ be complete noetherian local of dimension
$n>0$, and put $U=\Spec A-\{\mathfrak m\}$. Its dimension is at
most $n-1$, so induction represents $F_U$ by $X_U$.
For each point of the represented closed fibre, condition 5 gives
compatible finite local algebras $B_{i,r}$ over $A/\mathfrak m^r$,
with $B_{i,r+1}/\mathfrak m^rB_{i,r+1}=B_{i,r}$.
Their limit $B_i$ is finite over $A$. To verify this last statement,
lift a finite set of generators of $B_{i,1}$ compatibly through
the surjections. Successively correcting coefficients in powers
of $\mathfrak m$ expresses every compatible element as a linear
combination of those lifts with coefficients in the complete ring
$A$. This gives a surjection $A^d\to B_i$; its kernel is closed,
and the same approximation identifies $B_i/\mathfrak m^rB_i$
with $B_{i,r}$. Thus $B_i$ is a finite, complete local algebra.
Condition 4 gives the compatible element $\xi_i\in F(B_i)$.
The maps $Z_i=\Spec B_i\to F$ have the required Artinian
pro-modular universal property. Condition $7_{\mathrm{pro}}$
makes $(Z_i)_U\to X_U$ etale.

We check that these maps are open immersions with pairwise disjoint
images. For $i,j$, equality over $X_U$ is a closed relation in
$(Z_i\times_A Z_j)_U$, since $X_U$ is separated. If a point of this
relation existed off the diagonal, its closure in the finite
$A$-scheme $Z_i\times_A Z_j$ would meet the closed fibre. The
noetherian valuative specialization criterion gives a trait
through such a specialization (allowing a residue-field extension).
The two maps of this trait to $F$ agree generically, hence agree
by condition 6. On every Artinian quotient at its closed point,
the pro-modular universal property forces equality of the maps
to the closed-fibre representing components. Distinct indices are
therefore impossible; for the same index, it forces equality of
the two maps $B_i\to V/\mathfrak m_V^r$ for every $r$.
Krull intersection in the discrete valuation ring $V$ then makes
the two maps $B_i\to V$ equal, contradicting the off-diagonal
generic point. The etale relation has open-and-closed diagonal,
so absence of points in its complement proves equality as schemes,
not just as sets. Thus each etale map is a monomorphism and hence
an open immersion, and their images are disjoint.

Proposition~\ref{A:formal-gluing} glues these finite schemes to
$X_U$ and represents $F_A$. The result is locally quasi-finite by
construction and separated by condition 6 and the noetherian
valuative criterion. For a noncomplete noetherian local base,
Theorem~\ref{A:completion-descent} descends the representation
from its faithfully flat completion. This establishes actual
local representability in dimension $n$.
Theorem~\ref{A:local-models-suffice} now supplies the global step
in that dimension. Induction is complete. Every noetherian local
ring has finite dimension, so the same local-to-global theorem
also treats locally noetherian bases of unbounded dimension.

For necessity, let $F=h_X$ with $X\to S$ separated locally
quasi-finite. Over a complete local base, its local branch at a
closed-fibre point is finite and complete, by the henselian branch
lemma. The formal universal property identifies a given finite
pro-modular $B$ with this branch: their quotients at every finite
level agree, and both rings are complete. Hence the map from
$\Spec B$ is an open immersion and is locally modular everywhere.
The other necessary conditions are the usual sheaf, finite
presentation, local adic, finite infinitesimal, and separatedness
properties; corrected necessity of 8 was proved in
Proposition~\ref{A:8-literal}.
\end{proof}

The strengthened premise in this theorem is the premise explicitly used
on cote 131, scan 5 before the intended weakening. We have not deduced
it from the actual-modularity premise on scan 14. Accordingly this
corrected theorem and the now unconditional unramified corollary do not
by themselves settle the main literal theorem.
